%%%%

%%%   Самарский государственный технический университет

%%%   Кафедра прикладной математики и информатики

%%%

%%%   Преамбула для статьи в журнал "Вестник Самарского государственного

%%%   технического университета. Серия: «Физико-математические науки»"

%%%   Ориентирована под MikTEx 2.4 for Win32

%%%

%%%   Хотя сам журнал собирается под GNU/Linux на дистрибутиве TeXLive



\documentclass[11pt,a4paper,twoside]{article}



\usepackage[cp1251]{inputenc}

\usepackage[T2A]{fontenc}

\usepackage[english,russian]{babel}

\usepackage{samgtubib}

\usepackage[dvips]{graphicx}

\usepackage[multiple]{footmisc}



\usepackage{samgtu}

\renewcommand{\VestnikYear}{2009} % Год издания Вестника

\renewcommand{\VestnikNumber}{2\,(19)} % Номер Вестника

\renewcommand{\VestnikMonth}{Ноябрь} % Месяц выхода Вестника

\renewcommand{\VestnikData}{01.11.09} % Дата подписания Вестника в печать

\renewcommand{\VestnikUPL}{26,64} % Усл. печ. л.

\renewcommand{\VestnikUIL}{25,73} % Уч.-изд. л.

\renewcommand{\VestnikRegNumber}{479/09} % Регистрационный номер

\renewcommand{\VestnikOrder}{994} % Номер заказа









\thanks{}



%\begin{document}



\renewcommand{\firstpage}{102} %Начальная страница статьи для выноса в колонтитул

\renewcommand{\lastpage}{107} %Конечная страница статьи для выноса в колонтитул





\udk{621.52}



\title{Линеаризованная математическая модель синхронного электродвигателя при различных способах управления его скоростью}

      {Линеаризованная математическая модель синхронного электродвигателя  \dots}

\engtitle{Linearized mathematical model of synchronous motor with various methods of speed control} 

\author{М.\,С.~Лысов, А.\,В.~Стариков, В.\,А.~Стариков}

       {Л\,ы\,с\,о\,в~М.\,С., С\,т\,а\,р\,и\,к\,о\,в~А.\,В., С\,т\,а\,р\,и\,к\,о\,в~В.\,А. }

\engauthor{M.\,S.~Lysov, A.\,V.~Starikov, V.\,A.~Starikov}







\date{15.10.2007}





\address{\samgtu}

\engaddress{Samara State Technical University, Samara, Russia}



\email{mrcpk@rambler.ru}



\workabstract{Рассмотрены возможные способы управления скоростью синхронного электродвигателя. Приведены дифференциальные уравнения и построены структурные схемы синхронного электродвигателя как объекта при различных способах управления. Проведена линеаризация уравнений движения двигателя, получены передаточные функции по отношению к управляющим и возмущающему воздействиям.}



\engabstract{Various methods of synchronous motor speed control are considered. Differential equations and structures of synchronous motor being the object with various methods of control are presented. Equations of motor motion have been linearized, transfer functions related with control and perturbation have been obtained.}



\maketitle





На практике в основном применяют два способа регулирования скорости 

синхронного электродвигателя: изменением частоты питающего напряжения 

(тока), подаваемого на статорные обмотки; изменением амплитуды напряжения 

при коммутации статорных обмоток по датчику положения ротора (режим 

вентильного электродвигателя или бесколлекторного электродвигателя 

постоянного тока).





\begin{wrapfigure}{O}{.3\textwidth}

\centerline{\includegraphics{./00Pic/lysov1.eps}}

\caption{\small Модель обобщённого двухфазного синхронного электродвигателя}

\end{wrapfigure}



Для описания динамических режимов работы синхронного электродвигателя при 

любом способе управления его скоростью удобно обратиться к модели обобщённой 

двухфазной машины переменного тока (рис.~1)~[1].



Эта модель позволяет перейти к~условному векторному описанию электрических 

машин переменного тока. В рассмотрение вводятся две системы координат: 

$0\alpha \beta $ и~$0dq$. Система координат $0\alpha \beta $ неподвижна и~связана с~обмотками статора. 

Поскольку векторы удобно записывать в~комплексной форме, то введены действительная 

${\rm Re}$ и~мнимая ${\rm Im}$ оси, причём действительная ось направлена по оси $\alpha$. 

Система координат $0dq$ связана с~обмоткой ротора синхронной машины.



Суть модели обобщённой машины переменного тока заключается в~том, что 

действие обмоток статора $\alpha $ и~$\beta $, запитанных соответствующей системой 

переменного напряжения (системой токов $I_\alpha $ и~$I_\beta )$ и~создающих 

в~каждый конкретный момент времени векторы потокосцеплений $\psi _\alpha $ и~$\psi _\beta $, 

можно заменить действием одной обмотки, запитанной 

постоянным током $I_1 $ и~вращающейся со скоростью $\omega _0 $ магнитного 

поля.



В синхронном электродвигателе обмотка возбуждения, расположенная на роторе, 

запитывается постоянным напряжением. Она создаёт потокосцепление $\psi _{\text в} $ 

и~вращается вместе с~ротором со скоростью $\omega $.



Векторные уравнения для обмоток статора и~ротора синхронной машины запишутся 

следующим образом:

\begin{equation}

\label{lysov:eq1}

\bar {U}_1 = R_1 \bar {I}_1 + \frac{d\bar {\psi }_1 }{dt} + j\omega _0 \bar 

{\psi }_1 ,

\quad

\bar {U}_{\text в} = R_{\text в} \bar {I}_{\text в} + \frac{d\bar {\psi }_{\text в} }{dt} + j\omega \bar 

{\psi }_{\text в}.

\end{equation}

Здесь $R_1 $ и $R_{\text в} $ "--- активные сопротивления обмотки статора и~обмотки 

возбуждения; $U_1 $ и~$U_{\text в} $ "--- напряжения на статоре и~на обмотке 

возбуждения, соответственно. Недостаток системы (\ref{lysov:eq1}) заключается в~том, что 

первое и~второе уравнения записаны в~разных системах координат. Первое уравнение записано в~неподвижных координатах $0\alpha \beta $, второе "--- во вращающихся координатах $0dq$, что затрудняет их совместное решение. Перепишем (\ref{lysov:eq1}) в~единой 

системе координат $0xy$, вращающейся со скоростью магнитного поля $\omega _0$. 

При этом скорость вращения ротора относительно выбранных координат составит 

$\omega _0 - \omega $. Следовательно,

\begin{equation}

\label{lysov:eq2}

\bar {U}_1 = R_1 \bar {I}_1 + \frac{d\bar {\psi }_1 }{dt} + j\omega _0 \bar 

{\psi }_1 ,

\quad

\bar {U}_{\text в}  = R_{\text в}  \bar {I}_{\text в}  + \frac{d\bar {\psi }_{\text в}  }{dt} + j(\omega _0 - 

\omega )\bar {\psi }_{\text в}  .

\end{equation}





Потокосцепления статора и ротора создаются токами в обеих обмотках:

\begin{equation}

\label{lysov:eq3}

\bar {\psi }_1 = L_1 \bar {I}_1 + L_0 \bar {I}_{\text в}  ,

\quad

\bar {\psi }_{\text в}  = L_{\text в}  \bar {I}_{\text в}  + L_0 \bar {I}_1 ,

\end{equation}

где $L_1 $ и $L_{\text в}  $ "--- собственные индуктивности эквивалентной обмотки 

статора и~обмотки возбуждения; $L_0 $ "--- взаимная индуктивность. Выразим 

векторы токов в~(\ref{lysov:eq3}) через потокосцепления

\[

\bar {I}_1 = \frac{L_{\text в}  }{\Delta }\bar {\psi }_1 - \frac{L_0 }{\Delta }\bar 

{\psi }_{\text в}  ,

\quad

\bar {I}_{\text в}  = \frac{L_1 }{\Delta }\bar {\psi }_{\text в}  - \frac{L_0 }{\Delta }\bar 

{\psi }_1 

\]

и подставим их в~(\ref{lysov:eq2}):

\begin{equation}

\label{lysov:eq4}

\frac{d\bar {\psi }_1 }{dt} = \bar {U}_1 - \frac{R_1 L_{\text в}  }{\Delta }\bar 

{\psi }_1 + \frac{R_1 L_0 }{\Delta }\bar {\psi }_{\text в}  - j\omega _0 \bar {\psi 

}_1 ,

\quad

\frac{d\bar {\psi }_{\text в}  }{dt} = \bar {U}_{\text в}  - \frac{R_{\text в} L_1 }{\Delta }\bar 

{\psi }_{\text в}  + \frac{R_{\text в}  L_0 }{\Delta }\bar {\psi }_1 - j(\omega _0 - \omega 

)\bar {\psi }_{\text в}  ,

\end{equation}

где $\Delta = L_1 L_{\text в} - L_0^2 $.



Систему уравнений (\ref{lysov:eq4}) можно переписать в проекциях на оси $x$ и $y$:

\begin{equation}

\label{lysov:eq5}

\left\{

\begin{array}{l}

\displaystyle 

\frac{d\psi _{1x} }{dt} = U_{1x} - \frac{R_1 L_{\text в} }{\Delta }\psi _{1x} + 

\frac{R_1 L_0 }{\Delta }\psi _{\text вx} + \omega _0 \psi _{1y} ,\vspace{1mm}

\\

\displaystyle 

\frac{d\psi _{1y} }{dt} = U_{1y} - \frac{R_1 L_{\text в} }{\Delta }\psi _{1y} + 

\frac{R_1 L_0 }{\Delta }\psi _{\text вy} - \omega _0 \psi _{1x}; \vspace{1mm}

\\

\displaystyle 

\frac{d\psi _{\text вx} }{dt} = U_{\text вx} - \frac{R_{\text в} L_1 }{\Delta }\psi _{\text вx} + 

\frac{R_{\text в} L_0 }{\Delta }\psi _{1x} + (\omega _0 - \omega )\psi _{\text вy} ,\vspace{1mm}

\\

\displaystyle 

\frac{d\psi _{\text вy} }{dt} = U_{\text вy} - \frac{R_{\text в} L_1 }{\Delta }\psi _{\text вy} + 

\frac{R_{\text в} L_0 }{\Delta }\psi _{1y} - (\omega _0 - \omega )\psi _{\text вx} , 

\end{array}\right.

\end{equation}

где переменные с индексами $x$ и $y$ представляют собой проекции векторов на 

соответствующие оси.



С учётом того, что реальный трёхфазный синхронный электродвигатель можно 

свести с~помощью трёхфазно"=двухфазного преобразования к~обобщённой 

двухфазной модели, электромагнитный момент электродвигателя запишется так~[1]:

\begin{equation}

\label{lysov:eq6}

M = \frac{m_1 Z_{\text п} L_0 }{2\Delta }(\psi _{1y} \psi _{\text вx} - \psi _{1x} \psi_{\text вy} ).

\end{equation}

Здесь $m_1 $ "--- число фаз статора электродвигателя; $Z_{\text п} $ "--- число пар 

полюсов. Систему уравнений (\ref{lysov:eq5}) и~(\ref{lysov:eq6}) необходимо дополнить главным уравнением 

движения электропривода:

\begin{equation}

\label{lysov:eq7}

J\frac{d\omega }{dt} = M - M_{\text с},

\end{equation}

где $J$ "--- приведённый к~валу электродвигателя момент инерции; 

$M_{\text с} $ "--- момент сил сопротивления.



Переходя в (\ref{lysov:eq5})--(\ref{lysov:eq7}) к~операторной форме записи, после преобразований 

получим

\begin{equation}

\label{lysov:eq8}

\left\{

\begin{array}{l}

\displaystyle 

(T_1 p + 1)\psi _{1x} = T_1 U_{1x} + \frac{L_0 }{L_{\text в} }\psi _{\text вx} + T_1 

\omega _0 \psi _{1y} ,\vspace{1mm}

\\

\displaystyle 

(T_1 p + 1)\psi _{1y} = T_1 U_{1y} + \frac{L_0 }{L{\text в} }\psi _{\text вy} - T_1 

\omega _0 \psi _{1x} ;\vspace{1mm}

\\

\displaystyle 

(T_{\text в} p + 1)\psi _{\text вx} = T_{\text в} U_{\text вx} + \frac{L_0 }{L_1 }\psi _{1x} + T_{\text в} 

(\omega _0 - \omega )\psi _{\text вy} ,\vspace{1mm}

\\

\displaystyle 

(T_{\text в} p + 1)\psi _{\text вy} = T_{\text в} U_{\text вy} + \frac{L_0 }{L_1 }\psi _{1y} - T_{\text в} 

(\omega _0 - \omega )\psi _{\text вx};\vspace{1mm}

\\

\displaystyle 

Jp\omega = \frac{m_1 Z_{\text п} L_0 }{2\Delta }(\psi _{1y} \psi _{\text вx} - \psi _{1x} 

\psi _{\text вy} ) - M_{\text с }.

\end{array}\right.

\end{equation}

В (8) введены обозначения: $p$ "--- оператор дифференцирования; $T_1 = \frac{\Delta }{R_1 L_{\text в} }$, 

$T_{\text в} = \frac{\Delta }{R_{\text в} L_1 }$ "--- электромагнитные постоянные времени цепей статора и обмотки возбуждения, соответственно.



Система уравнений (\ref{lysov:eq8}) позволяет построить структурную схему синхронного 

электродвигателя как объекта управления (рис.~2). За входные управляющие 

воздействия приняты частота вращения магнитного поля $\omega _0 $ и~проекции 

векторов напряжений $U_{1x}$, $U_{1y}$, $U_{\text вx} $ и~$U_{\text вy} $. Основным 

возмущающим воздействием является момент $M_{\text c} $ сил сопротивления, 

а~выходной координатой "--- скорость $\omega $ вращения ротора электродвигателя.



Анализ уравнений (\ref{lysov:eq8}) и структурной схемы показывает, что синхронный 

электродвигатель представляет собой существенно нелинейный объект. Последнее 

обстоятельство объясняется, прежде всего, наличием в~структуре шести 

множительных звеньев.



\begin{figure}[b!]

\centerline{\includegraphics{./00Pic/lysov2.eps}}

\caption{Структурная схема синхронного электродвигателя как объекта  управления}

\end{figure}





Следует обратить внимание, что полученная структурная схема синхронного 

электродвигателя как объекта управления с~точностью до входных воздействий 

$U_{\text вx} $ и~$U_{\text вy} $ совпадает с~аналогичной схемой асинхронного двигателя~[2]. 

Поэтому можно провести линеаризацию системы уравнений (\ref{lysov:eq8}) методом 

разложения основных нелинейных зависимостей в~степенной ряд Тейлора в~окрестности 

некоторой рабочей точки с~параметрами: $\psi _{1x0} $, $\psi _{1y0} $, $\psi _{2x0} $ и $\psi _{2y0} $~[2]:

\begin{equation}

\label{lysov:eq9}

\left\{

\begin{array}{l}

\displaystyle

(T_1 p + 1)\psi _{1x} = T_1 U_{1x} + \frac{L_0 }{L_{\text в} }\psi _{\text вx} + T_1 \psi _{1y0} \omega _0, \vspace{1mm}

\\

\displaystyle

(T_1 p + 1)\psi _{1y} = T_1 U_{1y} + \frac{L_0 }{L_{\text в} }\psi _{{\text в}y} - T_1 \psi 

_{1x0} \omega _0;

\vspace{1mm}

\\

\displaystyle

(T_{\text в} p + 1)\psi _{{\text в}x} = U_{{\text в}x} + \frac{L_0 }{L_1 }\psi _{1x} + T_{\text в} \psi 

_{{\text в}y0} (\omega _0 - \omega ),

\vspace{1mm}

\\

\displaystyle

(T_{\text в} p + 1)\psi _{{\text в}x} = U_{{\text в}y} + \frac{L_0 }{L_1 }\psi _{1x} - T_{\text в} \psi 

_{{\text в}x0} (\omega _0 - \omega );

\vspace{1mm}

\\

\displaystyle

Jp\omega = \frac{mZ_{\text п} L_0 }{2\Delta }(\psi _{1y0} \psi _{{\text в}x} - \psi _{1x0} 

\psi _{{\text в}y} ) - M_{\text с}. 

\end{array}

\right.

\end{equation}





Система уравнений (\ref{lysov:eq9}) позволяет найти передаточные функции по отношению 

к~управляющим и~возмущающему воздействиям. Пользуясь принципом суперпозиции, 

найдём, прежде всего, передаточную функцию по отношению к~управляющему 

воздействию $\omega _0 $. Для этого приравняем нулю входные воздействия 

$U_{\text вx} $, $U_{\text вy} $ и~$M_{\text c} $. 



Передаточная функция синхронного электродвигателя по отношению 

к~управляющему воздействию $\omega _0 $ имеет такой же вид, как у~асинхронного 

электродвигателя~[2]:

\[

W_{\text{ду}1} (p) = \frac{\omega (p)}{\omega _0 (p)} = 

\frac{k_{\text{ду}1} \left( {\frac{T_1 }{k_{\text{ду}1} }p + 1} \right)}{T_{\text э} T_{\text м} T_1 p^3 + 

\frac{T_{\text э} T_{\text м} (T_1 + T_{\text в} )}{T_{\text в} }p^2 + 

\left[ {T_1 + \frac{T_{\text э} T_{\text м} }{T_{\text в} }\left( {1 - \frac{L_0^2 

}{L_1 L_{\text в} }} \right)} \right]p + 1},

\]

где коэффициент передачи двигателя по управлению

\[

k_{\text{ду}1} = 1 + \frac{T_1 (\psi _{1x0}^2 + \psi _{1y0}^2 )L_0 }{T_{\text в} 

(\psi _{1x0} \psi _{\text вx0} + \psi _{1y0} \psi _{\text вy0} )L_1 },

\]

а произведение приведённых электромагнитной и~электромеханической постоянных 

времени синхронного двигателя представляется так:

\[

T_{\text э} T_{\text м} = 

\frac{2\Delta J}{m_1 Z_{\text п} L_0 (\psi _{1x0} \psi _{\text вx0} + \psi_{1y0} \psi _{\text вy0} )}.

\]



Передаточная функция синхронного электродвигателя по отношению 

к~возмущающему воздействию $M_{\text с}$ имеет следующий вид:

\[

W_{\text{дв}1} (p) = \frac{\omega (p)}{M_{\text с} (p)} = 

- \frac{\displaystyle \frac{T_{\text э} T_{\text м} 

\left( {1 - \frac{L_0^2 }{L_1 L_{\text в} }} \right)}{JT_{\text в} }

\left( {\frac{T_1 T_{\text в} }{1 - \frac{L_0^2 }{L_1 L_{\text в} }}p^2 + 

\frac{T_1 + T_{\text в} }{1 - \frac{L_0^2 }{L_1 L_{\text в} }}p + 1} \right)}

{\displaystyle T_{\text э} T_{\text м} T_1 p^3 + \frac{T_{\text э} T_{\text м} (T_1 + T_{\text в} )}{T_{\text в} }p^2 + 

\left[ {T_1 + \frac{T_{\text э} T_{\text м} }{T_{\text в} }

\left( {1 - \frac{L_0^2 }{L_1 L_{\text в} }} \right)} \right]p + 1}.

\]



\begin{wrapfigure}{O}{.3\textwidth}

\centerline{\includegraphics{./00Pic/lysov3.eps}}

\caption{\small Модель двухфазного синхронного электродвигателя с~постоянными магнитами на роторе}

\vspace{-3mm}

\end{wrapfigure}



Таким образом, динамические характеристики синхронных машин, оснащённых 

обмоткой возбуждения, и~асинхронных двигателей <<в малом>> при частотном 

управлении совпадают.



В случае, когда синхронный электродвигатель имеет возбуждение от постоянных 

магнитов, расположенных на роторе (рис. 3), его математическая модель 

упрощается. 



Запишем первое уравнение в~(\ref{lysov:eq1}) в~системе координат $0dq$, вращающихся вместе 

с~ротором:

\begin{equation}

\label{lysov:eq10}

\bar {U}_1 = R_1 \bar {I}_1 + \frac{d\bar {\psi }_1 }{dt} + j(\omega - \omega _0 )\bar {\psi }.

\end{equation}

Здесь $\omega - \omega _0 $ "--- скорость вращения магнитного поля статора 

относительно выбранных координат. Запишем (\ref{lysov:eq10}) в~проекциях на оси $d$ и~$q$:

\begin{equation}

\label{lysov:eq11}

\left\{

\begin{array}{l}

 U_{1d} = R_1 I_{1d} + \frac{d\psi _{1d} }{dt} - (\omega - \omega _0 )\psi_{1q}, \\ 

 U_{1q} = R_1 I_{1q} + \frac{d\psi _{1q} }{dt} + (\omega - \omega _0 )\psi_{1d}.

 \end{array}

\right.

\end{equation}

Учтём, что в~рассматриваемом случае

\begin{equation}

\label{lysov:eq12}

\psi _{1d} = L_1 I_{1d} + \psi _{\text в} ,

\quad

\psi _{1q} = L_1 I_{1q}.

\end{equation}

Выражая из (\ref{lysov:eq12}) проекции вектора тока статора $I_{1d}$, $I_{1q} $ 

и~подставляя их в~(\ref{lysov:eq11}), получим систему уравнений цепи статора синхронного 

двигателя, выраженную через однородные переменные:

\begin{equation}

\label{lysov:eq13}

\left\{

\begin{array}{l}

 \displaystyle

U_{1d} = \frac{R_1 }{L_1 }(\psi _{1d} - \psi _{\text в} ) + \frac{d\psi _{1d} }{dt} - (\omega - \omega _0 )\psi _{1q},\vspace{1mm}\\

 \displaystyle 

U_{1q} = \frac{R_1 }{L_1 }\psi _{1q} + \frac{d\psi _{1q} }{dt} + (\omega - \omega _0 )\psi _{1d} .

\end{array}

\right.

\end{equation}



Момент, развиваемый синхронным электродвигателем [1], запишется так:

\begin{equation}

\label{lysov:eq14}

M = \frac{m_1 Z_{\text п} }{2}{\rm Im}\left\{ {\bar {\psi }_1^\ast \bar {I}_1 } \right\},

\end{equation}

где $\bar {\psi }_1^\ast = \psi _{1d} - j\psi _{1q} $ "--- сопряжённый вектор 

потокосцепления статора. Определяя мнимую часть произведения векторов в~(\ref{lysov:eq14}), получим

\begin{equation}

\label{lysov:eq15}

M = \frac{m_1 Z_{\text п} \psi _{\text в} }{2L_1 }\psi _{1q} .

\end{equation}



Уравнения (\ref{lysov:eq7}), (\ref{lysov:eq13}) и~(\ref{lysov:eq15}) дают полное описание синхронного электродвигателя 

в~системе частотного управления

\begin{equation}

\label{lysov:eq16}

\left\{

\begin{array}{l}

\displaystyle

\frac{d\psi _{1d} }{dt} = U_{1d} - \frac{R_1 }{L_1 }\psi _{1d} + \frac{R_1 }{L_1 }\psi_{\text в} + (\omega - \omega_0)\psi _{1q} ,

\vspace{1mm}

\\

\displaystyle

\frac{d\psi _{1q} }{dt} = U_{1q} - \frac{R_1 }{L_1 }\psi _{1q} - (\omega - \omega _0 )\psi _{1d} ,

\vspace{1mm}

\\

\displaystyle

J\frac{d\omega }{dt} = \frac{m_1 Z_{\text п} \psi _{\text в} }{2L_1 }\psi _{1q} - M_{\text с}. 

\end{array}

\right.

\end{equation}



Перейдём в~(\ref{lysov:eq16}) к операторной форме записи:

\begin{equation}

\label{lysov:eq17}

\left\{

\begin{array}{l}

\displaystyle 

(T_{11} p + 1)\psi _{1d} = T_{11} U_{1d} + \psi _{\text в} + T_{11} (\omega - \omega _0 )\psi _{1q},

\vspace{1mm}

\\ 

\displaystyle 

(T_{11} p + 1)\psi _{1q} = T_{11} U_{1q} - T_{11} (\omega - \omega _0 )\psi_{1d} ,

\vspace{1mm}

\\

\displaystyle 

J\frac{d\omega }{dt} = \frac{m_1 Z_{\text п} \psi _{\text в} }{2L_1 }\psi _{1q} - M_{\text с},

\end{array}

\right.

\end{equation}

где $T_{11} = \frac{L_1 }{R_1 }$.



\begin{wrapfigure}{O}{.45\textwidth}

\centerline{\includegraphics{./00Pic/lysov4.eps}}

\caption{\small Структурная схема синхронного электродвигателя с постоянными  магнитами}

\vspace{-6mm}

\end{wrapfigure}



Система уравнений (\ref{lysov:eq17}) позволяет построить структурную схему синхронного 

двигателя с~постоянными магнитами на роторе (рис.~4). 







Как и в предыдущем случае, синхронный электродвигатель представляет собой 

нелинейный объект управления, что вызвано наличием множительных звеньев. 

Проведём линеаризацию системы уравнений (\ref{lysov:eq17}) в~районе некоторой рабочей 

точки с~координатами $\psi _{1d0} $ и~$\psi _{1q0}$:

\begin{equation}

\label{lysov:eq18}

\left\{

\begin{array}{l}

 \displaystyle 

(T_{11} p + 1)\psi _{1d} = T_{11} U_{1d} + \psi_{\text в} + T_{11} (\omega - \omega _0 )\psi _{1q0} ,

\vspace{1mm}

\\

\displaystyle 

(T_{11} p + 1)\psi _{1q} = T_{11} U_{1q} - T_{11} (\omega - \omega _0 )\psi _{1d0} ,

\vspace{1mm}

\\

\displaystyle 

J\frac{d\omega }{dt} = \frac{m_1 Z_{\text п} \psi_{\text в} }{2L_1 }\psi _{1q} - M_{\text с}.

\end{array}

\right.\hspace{-1cm}

\end{equation}



Система уравнений (\ref{lysov:eq18}) позволяет найти передаточные функции синхронного 

электродвигателя с~постоянными магнитами на роторе как по управляющим 

воздействиям $\omega _0 $ и~$U_1 $, так и~по отношению к~основному 

возмущению $M_{\text с}$.



Передаточная функция синхронного двигателя по отношению к~изменению скорости 

вращения магнитного поля $\omega_0$ (частоты питающего напряжения) имеет вид

\[

W_{\text{ду}2} (p) = 

\frac{\omega (p)}{\omega _0 (p)} = \frac 1

{\frac{2JL_1 }{m_1 Z_{\text п} \psi _{1d0} \psi_{\text в} }p^2 + 

\frac{2JL_1 }{m_1 Z_{\text п} \psi _{1d0} \psi _{\text в} T_{11} }p + 1}

\]

или $$W_{\text{ду}2} (p) = \frac 1{T_{\text э1} T_{\text м1} p^2 + T_{\text м1} p + 1},$$

где $T_{\text э1} = T_{11} $, $T_{\text м1} = \frac{2JL_1 }{m_1 Z_{\text п} \psi _{1d0} \psi_{\text в} T_{11} }$.



\smallskip



Передаточная функция по отношению к возмущению $M_{\text с} $:

\[

W_{\text{ду}2} (p) = 

\frac{\omega (p)}{M_{\text с} (p)} = 

- \frac{ \frac{2L_1 }{m_1 Z_{\text п} \psi _{1d0} \psi_{\text в} T_{11} }

(T_{11} p + 1)}{ \frac{2JL_1 }{m_1 Z_{\text п} \psi _{1d0} \psi_{\text в} }p^2 + 

\frac{2JL_1 }{m_1 Z_{\text п} \psi _{1d0} \psi _{\text в} T_{11} }p + 1}.

\]



Все передаточные функции по управляющим воздействиям, полученные выше, 

описывают движение синхронного электродвигателя <<в малом>> при частотном 

способе управления скорости. 



В электроприводах с большим диапазоном регулирования скорости применяют 

способ управления синхронным электродвигателем по сигналам датчика положения 

ротора, привязанного к~магнитной системе машины~[3]. Как правило, в~этом 

случае синхронный электродвигатель оснащён постоянными магнитами, 

расположенными на роторе, а~управление скоростью производят изменением 

амплитуды (действующего значения) напряжения на обмотках статора. В~этом 

режиме в~синхронном двигателе поддерживают пространственный угол между 

векторами потокосцепления ротора и~статора, близким к~$90^\circ$, а~скорость 

магнитного поля $\omega_{0}$ практически совпадает со скоростью $\omega $ вращения ротора. 

Структурная схема, представленная на рис.~4, позволяет получить передаточную 

функцию синхронного электродвигателя и~в~этом случае.



Приравнивая нулю $M_{\text c} $ и $\omega _0 $, найдём передаточную функцию по 

отношению к~управляющему воздействию $U_{1q} $:

\begin{equation}

\label{lysov:eq19}

W_{\text{ду}3} (p) = \frac{\omega (p)}{U_{1q} (p)} = \frac{1}{\psi _{1d0} (T_{\text э1} 

T_{\text м1} p^2 + T_{\text м1} p + 1)}.

\end{equation}



Анализ передаточной функции (\ref{lysov:eq19}) показывает, что динамические характеристики 

синхронной машины, работающей в~режиме вентильного двигателя 

(бесколлекторного двигателя постоянного тока) аналогичны соответствующим 

характеристикам обычного электродвигателя постоянного тока. Скорость 

синхронного двигателя в~этом режиме пропорциональна напряжению, 

прикладываемому к~статорным обмоткам.









\begin{thebibliography}{3}



\item \emph{Михайлов, О.\,П.} Автоматизированный электропривод станков и промышленных роботов

[Текст]~/ О.\,П.~Михайлов. "--- М.: Машиностроение, 1990. "--- 304~с.





\item \emph{Стариков, А.\,В.} Линеаризованная математическая модель асинхронного 

электродвигателя как объекта системы частотного управления [Текст]~/ 

А.\,В.~Стариков~/\!/ Вестн. Сам. гос. техн. ун-та.  Сер.:

Физ.-мат. науки. "--- 2002. "--- \No~16. "--- С.~175--180. "--- ISBN 5--7964--0355--9.



\item \emph{Терехов, В.\,М.} Системы управления электроприводов: Учебник для студ. высш. 

учеб. заведений [Текст]~/ В.\,М. Терехов, О.\,И. Осипов; Под ред. В.\,М.~Терехова. "--- 

М.: Академик, 2005. "--- 304~с. 

\end{thebibliography}



\vspace{8mm}



\makeengtitle







%\end{document}

